\section{Valuative criteria}
\label{section:II.7}

In this section we give valuative criteria for separatedness and for
properness of a morphism, that is, criteria which involve a variable auxiliary
scheme of the form $\operatorname{Spec}(A)$, where $A$ is
a valuation ring. Under suitable ``Noetherian'' hypotheses, we can
restrict in these criteria to the case where $A$ is a \emph{discrete}
valuation ring. This will doubtless be the only case that we will
have to apply in what follows, and we have introduced arbitrary valuation
rings, in the general case, only in order to make the connection with
classical developments.

\subsection{Review of valuation rings.}
\label{subsection:II.7.1}

\begin{env}[7.1.1]
\label{II.7.1.1}
Among the various equivalent properties that characterize valuation
rings, we will use the following: a ring $A$ is said to be a
\emph{valuation ring} if it is an integral ring which is not a
field, and if, in the set of local rings contained in the field of
fractions $K$ of $A$ and distinct from $K$, $A$ is \emph{maximal} for the
relation of domination (\hyperref[I.8.1.1]{I, 8.1.1}). Recall that a
valuation ring is \emph{integrally closed}. If $A$ is a valuation ring,
$A_{\mathfrak{p}}$ is also a valuation ring for every prime ideal
$\mathfrak{p}\neq0$ of $A$.
\end{env}

\begin{env}[7.1.2]
\label{II.7.1.2}
Let $K$ be a field and $A$ a local subring of $K$ which is not a field;

\oldpage[II]{139}

then there exists a valuation ring having $K$ as field of fractions and
which dominates $A$ (\cite[p.~1-07, lemma~2]{I-1}).

On the other hand, let $B$ be a valuation ring, $k$ its residue field,
$K$ its field of fractions, and $L$ an extension of $k$. Then there exists a
\emph{complete} valuation ring $C$ which dominates $B$ and whose residue
field is $L$. Indeed, $L$ is an algebraic extension of a purely transcendental
extension $L'=k(T_\mu)_{\mu\in M}$; we know that we can extend the
valuation of $K$ corresponding to $B$ to a valuation of
$K'=K(T_\mu)_{\mu\in M}$ in such a way that $L'$ is the residue field of this
valuation (\cite[p.~98]{II-24}); replacing $B$ by the completion of the ring
of this extended valuation, we see that we can restrict to the case where $B$
is complete and $L$ is an algebraic closure of $k$. If $\overline{K}$ is an
algebraic closure of $K$, we can then extend to $\overline{K}$ the
valuation which defines $B$, and the corresponding residue field is an
algebraic closure of $k$, as we see by lifting to $\overline{K}$
the coefficients of a monic polynomial of $k[T]$. We are therefore finally
reduced to the case where $L=k$, and it then suffices to take for $C$ the
completion of $B$ to answer the question.
\end{env}

\begin{env}[7.1.3]
\label{II.7.1.3}
Let $K$ be a field and $A$ a subring of $K$; the integral closure $A'$
of $A$ in $K$ is the intersection of the valuation rings containing $A$ and
having $K$ as field of fractions
(\cite[p.~51, th.~2]{I-11}). Proposition
(\hyperref[II.7.1.2]{7.1.2}) can be expressed in the following equivalent
geometric form:
\end{env}

\begin{proposition}[7.1.4]
\label{II.7.1.4}
Let $Y$ be a prescheme, $p:X\to Y$ a morphism, $x$ a point of $X$,
$y=p(x)$, and $y'\neq y$ a specialization
(\hyperref[0.2.1.2]{0, 2.1.2}) of $y$. Then there exist a local scheme
$Y'$, the spectrum of a valuation ring, and a separated morphism $f:Y'\to Y$
such that, if $a$ denotes the unique closed point of $Y'$ and $b$ the generic
point of $Y'$, we have $f(a)=y'$ and $f(b)=y$. We may moreover suppose
that one of the following two additional properties is satisfied:
\begin{enumerate}
\item[\rm{(i)}] $Y'$ is the spectrum of a \emph{complete} valuation ring
whose residue field is algebraically closed, and there exists a
$k(y)$-homomorphism $k(x)\to k(b)$.
\item[\rm{(ii)}] There exists a $k(y)$-isomorphism
$k(x)\xrightarrow{\sim}k(b)$.
\end{enumerate}
\end{proposition}

Let $Y_1$ be the reduced closed sub-prescheme of $Y$ having
$\overline{\{y\}}$ as underlying space
(\hyperref[I.5.2.1]{I, 5.2.1}), and let $X_1$ be the closed sub-prescheme
that is the inverse image $p^{-1}(Y_1)$; since $y'\in\overline{\{y\}}$ by
hypothesis and $k(x)$ is the same in $X$ and in $X_1$, we may suppose $Y$
\emph{integral}, with generic point $y$; $\mathcal{O}_{y'}$ is then an
integral local ring which is not a field and whose field of fractions
is $\mathcal{O}_y=k(y)$, and $k(x)$ is an extension of $k(y)$. To
realize the conditions $f(a)=y'$ and $f(b)=y$ as well as the
additional condition (i) (resp. (ii)), we take
$Y'=\operatorname{Spec}(A')$, where $A'$ is a valuation ring dominating
$\mathcal{O}_{y'}$ which is complete and has an algebraically closed residue
field that is an extension of $k(x)$ (resp. a valuation ring $A'$
dominating $\mathcal{O}_{y'}$ and having $k(x)$ as its field of fractions);
the existence of $A'$ is guaranteed by
(\hyperref[II.7.1.2]{7.1.2}).

\begin{env}[7.1.5]
\label{II.7.1.5}
Recall that a local ring $A$ is said to be \emph{of dimension $1$} if there
exists a prime ideal distinct from the maximal ideal $\mathfrak{m}$, and if
every prime ideal of $A$ distinct from $\mathfrak{m}$ is a \emph{minimal}
prime ideal; when $A$ is \emph{integral}, it amounts to the same to say
that $\mathfrak{m}$ and $(0)$ are the only prime ideals and
$\mathfrak{m}\neq(0)$; in other words, $Y=\operatorname{Spec}(A)$
consists of two

\oldpage[II]{140}

points $a$, $b$: $a$ is the unique \emph{closed} point, we have
$\mathfrak{j}_a=\mathfrak{m}$, and $k(a)=k$ is the \emph{residue field}
$k=A/\mathfrak{m}$; $b$ is the \emph{generic point} of $Y$,
$\mathfrak{j}_b=(0)$, the set $\{b\}$ being the only open set of
$Y$ distinct from $\varnothing$ and from $Y$ (an open set which is therefore
\emph{everywhere dense}), and $k(b)=K$ is the \emph{field of fractions} of $A$.
\end{env}

\begin{env}[7.1.6]
\label{II.7.1.6}
For a local ring $A$ which is Noetherian and of dimension $1$, we know that
the following conditions are equivalent
(\cite[p.~2-08 and 17-01]{I-1}):
\begin{enumerate}
\item[\rm{a)}] $A$ is normal;
\item[\rm{b)}] $A$ is regular;
\item[\rm{c)}] $A$ is a valuation ring;
\end{enumerate}
moreover, $A$ is then a \emph{discrete valuation ring}. Propositions
(\hyperref[II.7.1.2]{7.1.2}) and
(\hyperref[II.7.1.3]{7.1.3}) then have the following analogues for
discrete valuation rings:
\end{env}

\begin{proposition}[7.1.7]
\label{II.7.1.7}
Let $A$ be a Noetherian integral local ring which is not a field, $K$ its
field of fractions, and $L$ an extension of finite type of $K$; then there
exists a discrete valuation ring having $L$ as field of fractions and
dominating $A$.
\end{proposition}

Suppose first that $L=K$. Let $\mathfrak{m}$ be the maximal ideal of $A$,
$(x_1,\ldots,x_n)$ a system of nonzero generators of $\mathfrak{m}$,
and $B$ the subring
$A[x_2/x_1,\ldots,x_n/x_1]$ of $K$, which is Noetherian. It is immediate that
the ideal $\mathfrak{m}B$ of $B$ is identical to the principal ideal $x_1B$;
if $\mathfrak{p}$ is a minimal prime ideal of $x_1B$, we know that
$\mathfrak{p}$ is of rank $1$ (\cite[t.~I, p.~277]{I-13}); in other words,
$B_{\mathfrak{p}}$ is a Noetherian local ring \emph{of dimension $1$};
it is clear that
$\mathfrak{p}B_{\mathfrak{p}}\cap A$ is an ideal of $A$ containing
$\mathfrak{m}$ and not containing $1$, hence equal to $\mathfrak{m}$, and
consequently $B_{\mathfrak{p}}$ \emph{dominates} $A$
(\hyperref[I.8.1.1]{I, 8.1.1}). It follows from the Krull--Akizuki theorem
(\cite[p.~293]{II-25}) that the integral closure $C$ of $B_{\mathfrak{p}}$
is a Noetherian ring (although $C$ is not necessarily a
$B_{\mathfrak{p}}$-module of finite type); if $\mathfrak{n}$ is a maximal
ideal of $C$, $C_{\mathfrak{n}}$ is a normal Noetherian local ring of
dimension $1$ (\cite[p.~295]{II-25}), hence a discrete valuation ring,
which dominates $B_{\mathfrak{p}}$ and \emph{a fortiori} $A$.

If now $L$ is an extension of finite type of $K$, we may, by the
foregoing, restrict to the case where $A$ is already a discrete
valuation ring. Let $w$ be a valuation of $K$ associated to $A$; there exists a
discrete valuation $w'$ of $L$ which \emph{extends} $w$: indeed, we reduce
by induction on the number of generators of $L$ to the case where
$L=K(\alpha)$, and then the proposition is classical
(\cite[p.~106]{II-24}).

\begin{corollary}[7.1.8]
\label{II.7.1.8}
Let $A$ be a Noetherian integral ring, $K$ its field of fractions, and $L$ an
extension of finite type of $K$. Then the integral closure of $A$ in $L$
is the intersection of the discrete valuation rings having $L$ as field of
fractions and containing $A$.
\end{corollary}

Indeed, such a discrete valuation ring, being normal, contains
\emph{a fortiori} every element of $L$ integral over $A$. It therefore
suffices to prove that, if $x\in L$ is not integral over $A$, there exists a
discrete valuation ring $C$ having $L$ as field of fractions, containing $A$
and not containing $x$. The hypothesis on $x$ indeed means that
$x\notin B=A[1/x]$; in other words, $1/x$ is not invertible in the
Noetherian ring $B$. There is therefore a prime ideal $\mathfrak{p}$ of $B$
containing $1/x$. The integral local ring $B_{\mathfrak{p}}$ is Noetherian and
contained in $L$, which is an extension of finite type of the field of
fractions of $B_{\mathfrak{p}}$ (the latter containing $K$). By virtue of
(\hyperref[II.7.1.7]{7.1.7}), there is therefore a discrete valuation ring
$C$ dominating $B_{\mathfrak{p}}$ and having $L$ as field of fractions;
since $1/x\in\mathfrak{p}B_{\mathfrak{p}}$ belongs to the maximal ideal of
$C$, we have $x\notin C$, which concludes the proof.

The geometric form of (\hyperref[II.7.1.7]{7.1.7}) is the following:

\oldpage[II]{141}
\begin{proposition}[7.1.9]
\label{II.7.1.9}
Let $Y$ be a locally Noetherian prescheme, $p:X\to Y$ a morphism
locally of finite type, $x$ a point of $X$, $y=p(x)$, and $y'\neq y$ a
specialization of $y$. Then there exist a local scheme $Y'$, the spectrum of a
discrete valuation ring, a separated morphism $f:Y'\to Y$ and a
rational $Y$-map $g$ from $Y'$ to $X$, such that, if $a$ denotes the
closed point of $Y'$ and $b$ its generic point, we have $f(a)=y'$,
$f(b)=y$, $g(b)=x$, and, in the commutative diagram
\[
\tag{7.1.9.1}
\label{II.7.1.9.1}
\begin{tikzcd}[column sep=large,row sep=large]
& k(x) \arrow[dl,"\gamma"'] \\
k(b) & k(y) \arrow[u,"\pi"] \arrow[l,"\varphi"']
\end{tikzcd}
\]
(where $\pi$, $\varphi$, $\gamma$ are the homomorphisms corresponding
to $p$, $f$ and $g$ respectively), $\gamma$ is a bijection.
\end{proposition}

As in (\hyperref[II.7.1.4]{7.1.4}), we can reduce to the case where $Y$
is integral with generic point $y$ (taking into account
(\hyperref[I.6.3.4]{I, 6.3.4, (iv)})), and, since the question is local on
$X$ and $Y$, we may suppose $p$ of finite type; we are then in the situation
of (\hyperref[II.7.1.4]{7.1.4}), with the additional property that
$k(x)$ is an extension \emph{of finite type} of $k(y)$
(\hyperref[I.6.4.11]{I, 6.4.11}) and that $\mathcal{O}_{y'}$ is Noetherian;
this makes it possible to apply (\hyperref[II.7.1.7]{7.1.7}) and to take
$Y'=\operatorname{Spec}(A')$, where $A'$ is a discrete valuation ring
dominating $\mathcal{O}_{y'}$ and having $k(x)$ as its field of fractions. We
have thus defined a commutative diagram
(\hyperref[II.7.1.9.1]{7.1.9.1}), where $\gamma$ is a bijection, and $\pi$
and $\varphi$ correspond to the morphisms $p$ and $f$. Moreover, since $X$ and
$Y$ are locally Noetherian (\hyperref[I.6.6.2]{I, 6.6.2}) and $Y'$
is integral, there is one and only one rational $Y$-map $g$ from $Y'$
to $X$ to which the isomorphism $\gamma$ corresponds
(\hyperref[I.7.1.15]{I, 7.1.15}), which completes the proof.

\subsection{Valuative criterion for separatedness}
\label{subsection:II.7.2}

\begin{proposition}[7.2.1]
\label{II.7.2.1}
Let $X$, $Y$ be two preschemes and $f:X\to Y$ a quasi-compact morphism. The
following two conditions are equivalent:
\begin{enumerate}
\item[\rm a)] The morphism $f$ is closed.
\item[\rm b)] For every $x\in X$ and every specialization $y'$ of $y=f(x)$
distinct from $y$, there exists a specialization $x'$ of $x$ such that
$f(x')=y'$.
\end{enumerate}
\end{proposition}

Condition b) expresses that $f(\overline{\{x\}})=\overline{\{y\}}$, and is
therefore a consequence of a). To show that b) implies a), consider a closed
subset $X'$ of the underlying space $X$; let $Y'=\overline{f(X')}$, and
let us show that $Y'=f(X')$. Consider the reduced closed sub-preschemes of
$X$ and $Y$ having $X'$ and $Y'$ respectively as underlying spaces
(\hyperref[I.5.2.1]{I, 5.2.1}); there is then a morphism $f':X'\to Y'$ such
that the diagram
\[
\begin{tikzcd}[column sep=large,row sep=large]
X' \arrow[r,"f'"] \arrow[d] & Y' \arrow[d] \\
X \arrow[r,"f"'] & Y
\end{tikzcd}
\]
is commutative (\hyperref[I.5.2.2]{I, 5.2.2}), and, since $f$ is
quasi-compact, so is $f'$. We are therefore reduced to proving that, if
$f$ is a quasi-compact and \emph{dominant} morphism, then
\oldpage[II]{142}
condition b) implies that $f(X)=Y$. Now let $y'$ be a point of $Y$, and let $y$
be the generic point of an irreducible component of $Y$ containing $y'$; by
virtue of b), it suffices to show that $f^{-1}(y)$ is not empty. But we know
that this property is a consequence of the fact that $f$ is dominant and
quasi-compact (\hyperref[I.6.6.5]{I, 6.6.5}).

\begin{corollary}[7.2.2]
\label{II.7.2.2}
Let $f:X\to Y$ be a quasi-compact immersion morphism. For the underlying
space $X$ to be closed in $Y$, it is necessary and sufficient that it contain
every specialization (in $Y$) of each of its points.
\end{corollary}

\begin{proposition}[7.2.3]
\label{II.7.2.3}
Let $Y$ be a prescheme (resp. a locally Noetherian prescheme) and $f:X\to Y$
a morphism (resp. a morphism locally of finite type). The following
conditions are equivalent:
\begin{enumerate}
\item[\rm a)] $f$ is separated.
\item[\rm b)] The diagonal morphism $X\to X\times_YX$ is quasi-compact, and,
for every $Y$-prescheme of the form $Y'=\operatorname{Spec}(A)$, where $A$ is a
valuation ring (resp. a discrete valuation ring), two
$Y$-morphisms from $Y'$ to $X$ which coincide at the generic point of $Y'$ are
equal.
\item[\rm c)] The diagonal morphism $X\to X\times_YX$ is quasi-compact, and,
for every $Y$-prescheme of the form $Y'=\operatorname{Spec}(A)$, where $A$ is a
valuation ring (resp. a discrete valuation ring), two
$Y'$-sections of $X'=X_{(Y')}$ which coincide at the generic point of $Y'$ are
equal.
\end{enumerate}
\end{proposition}

The equivalence of b) and c) follows from the bijective correspondence between
$Y$-morphisms from $Y'$ to $X$ and $Y'$-sections of $X'$
(\hyperref[I.3.3.14]{I, 3.3.14}). If $X$ is separated over $Y$, condition b)
is satisfied by virtue of (\hyperref[I.7.2.2.1]{I, 7.2.2.1}), since $Y'$
is integral. It remains to prove that b) implies that the diagonal morphism
$\Delta:X\to X\times_YX$ is closed, and it amounts to the same to show that it
satisfies the criterion of (\hyperref[II.7.2.2]{7.2.2}). Now let $z$ be a
point of the diagonal $\Delta(X)$ and $z'\neq z$ a specialization of $z$ in
$X\times_YX$. Then there exist (\hyperref[II.7.1.4]{7.1.4}) a
valuation ring $A$ and a morphism $f$ from $Y'=\operatorname{Spec}(A)$ to
$X\times_YX$ which maps the closed point $a$ of $Y'$ to $z'$ and the generic
point $b$ of $Y'$ to $z$; this morphism makes $Y'$ an
$(X\times_YX)$-prescheme, and \emph{a fortiori} a $Y$-prescheme. If we compose
the two projections of $X\times_YX$ with $f$, we obtain two $Y$-morphisms
$g_1$, $g_2$ from $Y'$ to $X$, which by hypothesis coincide at the point $b$;
they are therefore equal to one and the same morphism $g$, which means
(\hyperref[I.5.3.1]{I, 5.3.1}) that $f$ factors as $f=\Delta\circ g$, and
consequently $z'\in\Delta(X)$. When we suppose $Y$ locally Noetherian and $f$
locally of finite type, $X\times_YX$ is locally Noetherian
(\hyperref[I.6.6.7]{I, 6.6.7}); we can then repeat the preceding
reasoning, supposing in it that $A$ is a discrete valuation ring, by virtue
of (\hyperref[II.7.1.9]{7.1.9}).

\begin{remark}[7.2.4]
\label{II.7.2.4}
\begin{enumerate}
\item[\rm (i)] The hypothesis that the morphism $\Delta$ is quasi-compact is
always satisfied when $Y$ is locally Noetherian and $f$ locally of
finite type, since $X\times_YX$ is then locally Noetherian
(\hyperref[I.6.6.4]{I, 6.6.4, (i)}). In the general case, it also means
that, for every covering $(U_\alpha)$ of $X$ by affine open subsets, the
sets $U_\alpha\cap U_\beta$ are \emph{quasi-compact}.
\item[\rm (ii)] For $f$ to be separated, it \emph{suffices} that condition b)
or condition c) be satisfied for valuation rings $A$ which are
\emph{complete} and whose residue field is \emph{algebraically closed};
this indeed follows from the proof of
(\hyperref[II.7.2.3]{7.2.3}) and from
(\hyperref[II.7.1.4]{7.1.4}).
\end{enumerate}
\end{remark}

\oldpage[II]{143}
\subsection{Valuative criterion for properness}
\label{subsection:II.7.3}

\begin{proposition}[7.3.1]
\label{II.7.3.1}
Let $A$ be a valuation ring, $Y=\operatorname{Spec}(A)$, $b$ the generic
point of $Y$, $X$ an integral scheme, and $f:X\to Y$ a closed morphism such
that $f^{-1}(b)$ consists of a single point $x$ and the corresponding
homomorphism $k(b)\to k(x)$ is bijective. Then $f$ is an isomorphism.
\end{proposition}

Since $f$ is closed and dominant, we have $f(X)=Y$; it suffices
(\hyperref[I.4.2.2]{I, 4.2.2}) to prove that, for every $y'\neq b$ in $Y$,
there exists a \emph{single} point $x'$ such that $f(x')=y'$, and that the
corresponding homomorphism $\mathcal{O}_{y'}\to\mathcal{O}_{x'}$ is bijective,
since $f$ will then be a homeomorphism. Now, if $f(x')=y'$, $\mathcal{O}_{x'}$
is a local ring contained in $K=k(x)=k(b)$ and dominating $\mathcal{O}_{y'}$;
the latter is the local ring $A_{y'}$, hence is a valuation ring
(\hyperref[II.7.1.1]{7.1.1}) having $K$ as field of fractions. On the
other hand, we have $\mathcal{O}_{x'}\neq K$, since $x'$ is not the generic
point of $X$ (\hyperref[0.2.1.3]{0, 2.1.3}); we conclude that
$\mathcal{O}_{x'}=\mathcal{O}_{y'}$. Since $X$ is an integral scheme, the
relation $\mathcal{O}_{x'}=\mathcal{O}_{x''}$ implies $x'=x''$
(\hyperref[I.8.2.2]{I, 8.2.2}), which completes the proof.

\begin{env}[7.3.2]
\label{II.7.3.2}
Let $A$ be a valuation ring, $Y=\operatorname{Spec}(A)$, and $b$ the generic
point of $Y$, so that $\mathcal{O}_b=k(b)$ is equal to $K$, the field of
fractions of $A$; let $f:X\to Y$ be a morphism. We know
(\hyperref[I.7.1.4]{I, 7.1.4}) that the \emph{rational} $Y$-sections of
$X$ are in bijective correspondence with the \emph{germs} of $Y$-sections
(defined in neighbourhoods of $b$) at the point $b$, whence a canonical
map
\[
\tag{7.3.2.1}
\label{II.7.3.2.1}
\Gamma_{\mathrm{rat}}(X/Y)\longrightarrow
\Gamma(f^{-1}(b)/\operatorname{Spec}(K))
\]
the elements of $\Gamma(f^{-1}(b)/\operatorname{Spec}(K))$ being moreover
identified, by definition (\hyperref[I.3.4.5]{I, 3.4.5}), with the
\emph{points of $f^{-1}(b)=X\otimes_AK$ rational over $K$}. When $f$ is
\emph{separated}, it follows from (\hyperref[I.5.4.7]{I, 5.4.7}) that the
map (\hyperref[II.7.3.2.1]{7.3.2.1}) is \emph{injective}, since $Y$ is
an integral scheme.

Composing (\hyperref[II.7.3.2.1]{7.3.2.1}) with the canonical map
$\Gamma(X/Y)\to\Gamma_{\mathrm{rat}}(X/Y)$
(\hyperref[I.7.1.2]{I, 7.1.2}), we obtain a canonical map
\[
\tag{7.3.2.2}
\label{II.7.3.2.2}
\Gamma(X/Y)\longrightarrow \Gamma(f^{-1}(b)/\operatorname{Spec}(K)).
\]
When $f$ is \emph{separated}, this map is again \emph{injective}
(\hyperref[I.5.4.7]{I, 5.4.7}).
\end{env}

\begin{proposition}[7.3.3]
\label{II.7.3.3}
Let $A$ be a valuation ring with field of fractions $K$,
$Y=\operatorname{Spec}(A)$, $b$ the generic point of $Y$, and $f:X\to Y$ a
\emph{separated} and \emph{closed} morphism. Then the canonical map
(\hyperref[II.7.3.2.2]{7.3.2.2}) is bijective (which amounts to saying
that it is \emph{surjective}, and implies that the rational $Y$-sections of
$X$ are \emph{everywhere defined}).
\end{proposition}

So let $x$ be a point of $f^{-1}(b)$ \emph{rational over $K$}. Since $f$ is
separated, so is the morphism
$f^{-1}(b)\to\operatorname{Spec}(K)$ corresponding to $f$
(\hyperref[I.5.5.1]{I, 5.5.1, (iv)}), and, every section of $f^{-1}(b)$
being a closed immersion (\hyperref[I.5.4.6]{I, 5.4.6}), $\{x\}$ is closed
\emph{in $f^{-1}(b)$}. Consider the reduced closed sub-prescheme $X'$ of
$X$ having as underlying space the closure $\overline{\{x\}}$ of $\{x\}$
in $X$. It is clear that the restriction of $f$ to $X'$ satisfies the hypotheses
of (\hyperref[II.7.3.1]{7.3.1}), hence is an \emph{isomorphism} of $X'$ onto
$Y$, whose inverse isomorphism is the desired $Y$-section of $X$.

\begin{env}[7.3.4]
\label{II.7.3.4}
To state the following two results, we will use a
terminology which will be justified and discussed in Chapter IV: if $F$ is
a subset
\oldpage[II]{144}
of a prescheme $Y$, we call \emph{codimension} of $F$ in $Y$, and denote by
$\operatorname{codim}_YF$, the infimum of the integers
$\dim(\mathcal{O}_z)$, where $z$ runs through $F$.
\end{env}

\begin{corollary}[7.3.5]
\label{II.7.3.5}
Let $Y$ be a reduced locally Noetherian prescheme and $N$ the set of
points $y\in Y$ at which $Y$ is not regular
(\hyperref[0.4.1.4]{0, 4.1.4}); suppose that
$\operatorname{codim}_YN\geq2$. Let $f:X\to Y$ be a separated and closed
morphism of finite type, and let $g$ be a rational $Y$-section of $X$; if $Y'$
is the set of points of $Y$ at which $g$ is not defined (a set which is
\emph{closed} (\hyperref[I.7.2.1]{I, 7.2.1})), we have
$\operatorname{codim}_YY'\geq2$.
\end{corollary}

It suffices to prove that $g$ is defined at every point $z\in Y$ such that
$\dim\mathcal{O}_z\leq1$. If $\dim\mathcal{O}_z=0$, $z$ is the generic point
of an irreducible component of $Y$ (\hyperref[I.1.1.14]{I, 1.1.14}), hence
belongs to every open subset of $Y$ that is everywhere dense, and in particular
to the domain of definition of $g$. Suppose therefore that $\dim\mathcal{O}_z=1$;
$\mathcal{O}_z$ is then by hypothesis a regular Noetherian local ring,
hence (\hyperref[II.7.1.6]{7.1.6}) a \emph{discrete valuation ring}.
Let $Z=\operatorname{Spec}(\mathcal{O}_z)$; since $U=Y-Y'$ is everywhere
dense, $U\cap Z$ is not empty (\hyperref[I.2.4.2]{I, 2.4.2}); let $g'$
be the rational map from $Z$ to $X$ induced by $g$
(\hyperref[I.7.2.8]{I, 7.2.8}); it suffices to show that $g'$ is a
\emph{morphism} (\hyperref[I.7.2.9]{I, 7.2.9}). Now $g'$ can be
regarded as a rational $Z$-section of the $Z$-prescheme
$f^{-1}(Z)=X\times_YZ$; it is clear that the morphism $f^{-1}(Z)\to Z$
corresponding to $f$ is closed, and it follows from
(\hyperref[I.5.5.1]{I, 5.5.1, (i)}) that it is separated; we then conclude from
(\hyperref[II.7.3.3]{7.3.3}) that $g'$ is everywhere defined; since $Z$ is
reduced and $X$ is separated over $Y$, $g'$ is a morphism
(\hyperref[I.7.2.2]{I, 7.2.2}).

\begin{corollary}[7.3.6]
\label{II.7.3.6}
Let $S$ be a locally Noetherian prescheme and $X$ and $Y$ two
$S$-preschemes; suppose $Y$ reduced, and moreover that the set $N$ of
points $y\in Y$ at which $Y$ is not regular is such that
$\operatorname{codim}_YN\geq2$; suppose finally that the structural morphism
$X\to S$ is proper. Let $f$ be a rational $S$-map from $Y$ to
$X$, and let $Y'$ be the set of points of $Y$ at which $f$ is not defined;
then we have $\operatorname{codim}_YY'\geq2$.
\end{corollary}

We know (\hyperref[I.7.1.2]{I, 7.1.2}) that we can identify the
rational $S$-maps from $Y$ to $X$ with the rational $Y$-sections
of $X\times_SY$; since the structural morphism $X\times_SY\to Y$ is closed
(\hyperref[II.5.4.1]{5.4.1}), we can apply
(\hyperref[II.7.3.5]{7.3.5}), whence the corollary.

\begin{remark}[7.3.7]
\label{II.7.3.7}
The hypotheses made on $Y$ in (\hyperref[II.7.3.5]{7.3.5}) and
(\hyperref[II.7.3.6]{7.3.6}) will in particular be fulfilled when $Y$
is \emph{normal} (\hyperref[0.4.1.4]{0, 4.1.4}), by virtue of
(\hyperref[II.7.1.6]{7.1.6}).
\end{remark}

We can characterize universally closed (resp. proper) morphisms
by a converse of (\hyperref[II.7.3.3]{7.3.3}):

\begin{theorem}[7.3.8]
\label{II.7.3.8}
Let $Y$ be a prescheme (resp. a locally Noetherian prescheme) and $f:X\to Y$
a quasi-compact separated morphism (resp. a morphism of finite type). The
following conditions are equivalent:
\begin{enumerate}
\item[\rm a)] $f$ is universally closed (resp. proper).
\item[\rm b)] For every $Y$-scheme of the form
$Y'=\operatorname{Spec}(A)$, where $A$ is a valuation ring (resp. a discrete
valuation ring) with field of fractions $K$, the canonical map
\[
\operatorname{Hom}_Y(Y',X)\longrightarrow
\operatorname{Hom}_Y(\operatorname{Spec}(K),X)
\]
corresponding to the canonical injection $A\to K$ is surjective (resp.
bijective).
\oldpage[II]{145}
\item[\rm c)] For every $Y$-scheme of the form $Y'=\operatorname{Spec}(A)$,
where $A$ is a valuation ring (resp. a discrete valuation ring),
the canonical map (\hyperref[II.7.3.2.2]{7.3.2.2}) relative to the
$Y'$-prescheme $X_{(Y')}$ is surjective (resp. bijective).
\end{enumerate}
\end{theorem}

The equivalence of b) and c) follows immediately from
(\hyperref[I.3.3.14]{I, 3.3.14}); a) implies c), because a) implies, under
either hypothesis, that $f_{(Y')}$ is separated
(\hyperref[I.5.5.1]{I, 5.5.1, (iv)}) and closed, and it suffices to apply
(\hyperref[II.7.3.3]{7.3.3}). It remains to prove that b) implies a).
Let us first place ourselves in the case where $Y$ is arbitrary and $f$
separated and quasi-compact. If condition b) is satisfied by $f$, it is also
satisfied by $f_{(Y'')}:X_{(Y'')}\to Y''$, where $Y''$ is an arbitrary
$Y$-prescheme, by virtue of the equivalence of b) and c) and of the fact that
$X_{(Y'')}\times_{Y''}Y'=X\times_YY'$ for every morphism $Y'\to Y''$
(\hyperref[I.3.3.9.1]{I, 3.3.9.1}); since moreover $f_{(Y'')}$ is
separated and quasi-compact when $f$ is
(\hyperref[I.5.5.1]{I, 5.5.1, (iv)} and
\hyperref[I.6.6.4]{6.6.4, (iii)}), we are reduced to proving that b) implies
that $f$ is \emph{closed}. For this, it suffices to verify condition b) of
(\hyperref[II.7.2.1]{7.2.1}). So let $x\in X$, and let $y'$ be a
specialization of $y=f(x)$ distinct from $y$; by virtue of
(\hyperref[II.7.1.4]{7.1.4}), there are a scheme $Y'$, the spectrum of a
valuation ring, and a separated morphism $g:Y'\to Y$ such that, if $a$
denotes the closed point and $b$ the generic point of $Y'$, we have
$g(a)=y'$, $g(b)=y$, and there exists a $k(y)$-homomorphism
$k(x)\to k(b)$. The latter corresponds canonically to a $Y$-morphism
$\operatorname{Spec}(k(b))\to X$
(\hyperref[I.2.4.6]{I, 2.4.6}), and it therefore follows from b) that there
exists a $Y$-morphism $h:Y'\to X$ to which the preceding morphism
corresponds. We then have $h(b)=x$; if we set $h(a)=x'$, $x'$ is a
specialization of $x$, and we have $f(x')=f(h(a))=g(a)=y'$.

If now $Y$ is locally Noetherian and $f$ of finite type, hypothesis b)
first implies that $f$ is \emph{separated}, by virtue of
(\hyperref[II.7.2.3]{7.2.3}), the diagonal morphism $X\to X\times_YX$
being quasi-compact (\hyperref[II.7.2.4]{7.2.4}). Moreover, to verify
that $f$ is proper, it suffices to show that $f_{(Y'')}:X_{(Y'')}\to Y''$ is
\emph{closed} for every $Y$-prescheme $Y''$ \emph{of finite type}, taking into
account (\hyperref[II.5.6.3]{5.6.3}). Since $Y''$ is then locally
Noetherian, we can repeat the reasoning made in the first case, taking for
$Y'$ the spectrum of a discrete valuation ring and applying
(\hyperref[II.7.1.9]{7.1.9}) instead of
(\hyperref[II.7.1.4]{7.1.4}).

\begin{remark}[7.3.9]
\label{II.7.3.9}
\begin{enumerate}
\item[\rm (i)] When $Y$ is an arbitrary prescheme and $f$ a separated
morphism, for $f$ to be universally closed it \emph{suffices} that
condition b) or condition c) be satisfied for the valuation rings
$A$ which are \emph{complete} and whose residue field is \emph{algebraically
closed}; this indeed follows from the proof and from
(\hyperref[II.7.1.4]{7.1.4}).
\item[\rm (ii)] From criterion c) of
(\hyperref[II.7.3.8]{7.3.8}) we deduce a new proof of the fact that a
projective morphism $X\to Y$ is closed (\hyperref[II.5.5.3]{5.5.3}), closer
to the classical methods. Indeed, we may suppose $Y$ affine, and consequently
$X$ identified with a closed sub-prescheme of a projective bundle
$\mathbf{P}_Y^n$ (\hyperref[II.5.3.3]{5.3.3}); to prove that $X\to Y$
is closed, it suffices to verify that this is so for the structural morphism
$\mathbf{P}_Y^n\to Y$, and criterion c) of
(\hyperref[II.7.3.8]{7.3.8}), together with
(\hyperref[II.4.1.3.1]{4.1.3.1}), shows that we are reduced to proving the
following fact: \emph{if $Y$ is the spectrum of a valuation ring $A$ with
field of fractions $K$, every point of $\mathbf{P}_Y^n$ with values in $K$
comes (by restriction to the generic point of $Y$) from a point of
$\mathbf{P}_Y^n$ with values in $A$}. Now every invertible
$\mathcal{O}_Y$-module is trivial (\hyperref[I.2.4.8]{I, 2.4.8}); hence it
follows from (\hyperref[II.4.2.6]{4.2.6}) that a point of $\mathbf{P}_Y^n$
with values in $K$ is identified with a class of elements $(\zeta c_0,\zeta c_1,\ldots,\zeta
c_n)$ of $K$, where $\zeta\neq0$ and the $c_i$ are elements of $K$ that are
not all zero. Now, by multiplying the $c_i$ by an element of $A$ of
\oldpage[II]{146}
suitable valuation, we may suppose that the $c_i$ all belong to
$A$ and that at least one of them is invertible. But then
(\hyperref[II.4.2.6]{4.2.6}) the system $(c_0,\ldots,c_n)$ also defines
a point of $\mathbf{P}_Y^n$ with values in $A$, which proves our
assertion.
\item[\rm (iii)] The criteria (\hyperref[II.7.2.3]{7.2.3}) and
(\hyperref[II.7.3.8]{7.3.8}) are especially convenient when we regard the
datum of a $Y$-prescheme $X$ as equivalent to the datum of the functor
\[
X(Y')=\operatorname{Hom}_Y(Y',X)
\]
in a $Y$-prescheme $Y'$; these criteria will allow us, for example, to
prove that, under certain conditions, the ``Picard schemes'' are proper.
\end{enumerate}
\end{remark}

\begin{corollary}[7.3.10]
\label{II.7.3.10}
Let $Y$ be an integral scheme (resp. a locally Noetherian integral scheme),
$X$ an integral scheme, and $f:X\to Y$ a dominant morphism.
\begin{enumerate}
\item[\rm (i)] If $f$ is quasi-compact and universally closed, every valuation
ring whose field of fractions is the field $R(X)$ of rational functions
on $X$ and which dominates a local ring of $Y$ also dominates a
local ring of $X$.
\item[\rm (ii)] Conversely, suppose that $f$ is of finite type and that the
property stated in (i) is satisfied by every valuation ring (resp.
every discrete valuation ring) having $R(X)$ as field of fractions.
Then $f$ is proper.
\end{enumerate}
\end{corollary}

Let us first note that the hypotheses imply in all cases that $f$ is
separated (\hyperref[I.5.5.9]{I, 5.5.9}).
\begin{enumerate}
\item[\rm (i)] Let $K=R(Y)$, $L=R(X)$, $y$ a point of $Y$, and $A$ a valuation
ring having $L$ as field of fractions and dominating
$\mathcal{O}_y$; the injection $\mathcal{O}_y\to A$ therefore defines a morphism
$h$ from $Y'=\operatorname{Spec}(A)$ to $Y$
(\hyperref[I.2.4.4]{I, 2.4.4}) such that $h(a)=y$, denoting by $a$ the
closed point of $Y'$; moreover, if $\eta$ is the generic point of $Y$, which
is also that of $\operatorname{Spec}(\mathcal{O}_y)$, we have $h(b)=\eta$,
denoting by $b$ the generic point of $Y'$ (since $K\subset L$ by
hypothesis). If $\xi$ is the generic point of $X$, we have $k(\xi)=k(b)=L$ by
hypothesis, whence a $Y$-morphism $g:\operatorname{Spec}(L)\to X$ such that
$g(b)=\xi$; by virtue of (\hyperref[II.7.3.8]{7.3.8}), $g$ comes from a
$Y$-morphism $g':Y'\to X$. If $x=g'(a)$, it is clear that $A$ dominates
$\mathcal{O}_x$.
\item[\rm (ii)] The question being local on $Y$, we may always suppose
that $Y$ is affine (resp. affine and Noetherian). Since $f$ is of finite type,
we can apply Chow's lemma (\hyperref[II.5.6.1]{5.6.1}) in both cases.
There are therefore a projective morphism
$p:P\to Y$, an immersion morphism $j:X'\to P$ and a projective,
surjective and birational morphism $g:X'\to X$ ($X'$ being integral) such that
the diagram
\[
\begin{tikzcd}[column sep=large,row sep=large]
P \arrow[d,"p"'] & X' \arrow[l,"j"'] \arrow[d,"g"] \\
Y & X \arrow[l,"f"]
\end{tikzcd}
\]
is commutative. It suffices to prove that $j$ is a \emph{closed} immersion,
since then $f\circ g=p\circ j$ will be a projective morphism, hence proper, and,
since $g$ is surjective, $f$ will also be proper
(\hyperref[II.5.4.3]{5.4.3}). Let $Z$ be the reduced closed sub-prescheme of
$P$ having $\overline{j(X')}$ as underlying space
(\hyperref[I.5.2.1]{I, 5.2.1}); since $X'$ is integral, $j$ factors
as $i\circ h$, where $i:Z\to P$ is the canonical injection and $h:X'\to Z$ a
dominant open immersion (\hyperref[I.5.2.3]{I, 5.2.3}), and $Z$ is
integral;

\oldpage[II]{147}
moreover $Z$ is projective over $Y$, and we see that we may restrict to the
case where $P$ is \emph{integral} and $j$ is \emph{dominant and birational},
and everything comes down to seeing that $j$ is \emph{surjective}. Now let
$z\in P$; $\mathcal{O}_z$ is an integral local ring (resp. an integral and
Noetherian local ring) whose field of fractions is
\[
L=R(P)=R(X')=R(X).
\]
We may restrict to the case where $z$ is not the generic point of $P$.
There is consequently (\hyperref[II.7.1.2]{7.1.2} and
\hyperref[II.7.1.7]{7.1.7}) a valuation ring (resp. a discrete
valuation ring) $A$ having $L$ as field of fractions and dominating
$\mathcal{O}_z$. \emph{A fortiori}, $A$ dominates $\mathcal{O}_y$, setting
$y=p(z)$, and by hypothesis there is therefore an $x\in X$ such that $A$
dominates $\mathcal{O}_x$. Since $g$ is proper, the first part of the proof
shows that $A$ also dominates some $\mathcal{O}_{x'}$, for an $x'\in X'$; it
follows that $\mathcal{O}_z$ and $\mathcal{O}_{j(x')}=\mathcal{O}_{x'}$ are
related (\hyperref[I.8.1.4]{I, 8.1.4}), and, since $P$ is a scheme, this
implies $z=j(x')$ (\hyperref[I.8.2.2]{I, 8.2.2}) and completes the proof.
\end{enumerate}

\begin{corollary}[7.3.11]
\label{II.7.3.11}
Let $X$, $Y$ be two integral schemes and $f:X\to Y$ a dominant,
quasi-compact and universally closed morphism. Suppose moreover that $Y$ is
affine with (integral) ring $B$. Then $\Gamma(X,\mathcal{O}_X)$ is canonically
isomorphic to a subring of the integral closure of $B$ in $R(X)$.
\end{corollary}

Indeed (\hyperref[I.8.2.1.1]{I, 8.2.1.1}),
$\Gamma(X,\mathcal{O}_X)$ is identified with the intersection of the
$\mathcal{O}_x$ for $x\in X$; by virtue of (\hyperref[II.7.3.10]{7.3.10}),
(\hyperref[II.7.1.2]{7.1.2}) and (\hyperref[II.7.1.3]{7.1.3}),
$\Gamma(X,\mathcal{O}_X)$ is consequently contained in the intersection of the
valuation rings containing $B$ and having $R(X)$ as field of fractions;
the conclusion then follows from (\hyperref[II.7.1.3]{7.1.3}).

\begin{remark}[7.3.12]
\label{II.7.3.12}
Under the hypotheses of (\hyperref[II.7.3.11]{7.3.11}), and when we suppose
that $R(X)$ is an extension of finite type of $R(Y)$, we will in many
cases be able to conclude that $\Gamma(X,\mathcal{O}_X)$ is a \emph{module of
finite type} over the ring $B=\Gamma(Y,\mathcal{O}_Y)$. This will be the case,
for example, when $B$ is an \emph{algebra of finite type over a field}, since
we then know that the integral closure of $B$ in an extension of finite type
of its field of fractions is a $B$-module of finite type ([13], t.~I, p.~267,
th.~9); the conclusion then follows from (\hyperref[II.7.3.11]{7.3.11}) and
from the fact that $B$ is Noetherian.

In particular, \emph{a scheme $X$ that is proper and affine over a field $K$ is
finite}. Indeed, by virtue of (\hyperref[II.1.6.4]{1.6.4}),
(\hyperref[II.5.4.6]{5.4.6}) and
(\hyperref[I.6.4.4]{I, 6.4.4, c)}), we may restrict to the case where $X$ is
\emph{reduced}. Moreover, it will suffice to prove that each of the closed
sub-preschemes of $X$ having as underlying space an irreducible component of $X$
(these being finite in number) is finite over $K$, so that (taking into
account (\hyperref[II.5.4.5]{5.4.5})) we are finally reduced to the case
where $X$ is \emph{integral}. But then the result follows from the remarks made
above.

In Chapter III, we will recover this last proposition by other
methods and as a consequence of more general results, showing that, if
$f:X\to Y$ is proper and $Y$ locally Noetherian, $f_*(\mathcal{F})$ is
coherent for every coherent $\mathcal{O}_X$-module $\mathcal{F}$
(\hyperref[III.4.4.2]{III, 4.4.2}).

Let us note finally that the criterion (\hyperref[II.7.3.10]{7.3.10}) is
taken as the \emph{definition} of proper morphisms in classical Algebraic
Geometry. We have mentioned it only for this reason, the criterion
(\hyperref[II.7.3.8]{7.3.8}) seeming more manageable in all the
applications known to us.
\end{remark}

\oldpage[II]{148}


\subsection{Algebraic curves and function fields of dimension one}
\label{subsection:II.7.4}

The aim of this subsection is to show how the classical notion of algebraic
curve (as presented, for example, in the book of C.~Chevalley [23]) is
formulated in the language of schemes. Throughout this subsection, \emph{we
denote by $k$ a field, all the schemes considered are $k$-schemes of
finite type, and all morphisms are $k$-morphisms}.

\begin{proposition}[7.4.1]
\label{II.7.4.1}
Let $X$ be a prescheme of finite type over $k$ (hence Noetherian); let $x_i$
$(1\leq i\leq n)$ be the generic points of the irreducible components $X_i$ of
$X$, and let $K_i=k(x_i)$ $(1\leq i\leq n)$. The following conditions are
equivalent:
\begin{enumerate}
\item[\rm a)] Each of the $K_i$ is an extension of $k$ whose transcendence
degree is equal to 1.
\item[\rm b)] For every closed point $x$ of $X$, the local ring
$\mathcal{O}_x$ is of dimension 1 (\hyperref[II.7.1.5]{7.1.5}).
\item[\rm c)] The irreducible closed subsets of $X$ distinct from the $X_i$
are the closed points of $X$.
\end{enumerate}
\end{proposition}

Since $X$ is quasi-compact, every irreducible closed subset $F$ of $X$
contains a closed point (\hyperref[0.2.1.3]{0, 2.1.3}). By virtue of
(\hyperref[I.2.4.2]{I, 2.4.2}), there is a bijective correspondence between the
prime ideals of $\mathcal{O}_x$ and the irreducible closed subsets of $X$
containing $x$ (\hyperref[I.1.1.14]{I, 1.1.14}); the equivalence of b) and c)
follows immediately. On the other hand, if $\mathfrak{p}_\alpha$
$(1\leq\alpha\leq r)$ are the minimal prime ideals of the Noetherian local
ring $\mathcal{O}_x$, the local rings
$\mathcal{O}_x/\mathfrak{p}_\alpha$ are integral and have as fields of
fractions those $K_i$ such that $x\in X_i$. Moreover, we know ([1],
p.~4-06, th.~2) that the dimension of an integral local $k$-algebra of finite
type is equal to the transcendence degree over $k$ of its field of fractions.
Finally, the dimension of $\mathcal{O}_x$ is the supremum of the dimensions of the
$\mathcal{O}_x/\mathfrak{p}_\alpha$; now condition a) implies that these
dimensions are equal to 1, so a) implies b); conversely, if
$\mathcal{O}_x$ is of dimension 1, none of the $\mathfrak{p}_\alpha$ can
be equal to the maximal ideal of $\mathcal{O}_x$, since otherwise $\mathcal{O}_x$
would be of dimension 0; hence each of the
$\mathcal{O}_x/\mathfrak{p}_\alpha$ is of dimension 1, which shows that b)
implies a).

Note that, under the conditions of (\hyperref[II.7.4.1]{7.4.1}),
the set $X$ is \emph{empty or infinite}, as follows immediately from
(\hyperref[I.6.4.4]{I, 6.4.4}).

\begin{definition}[7.4.2]
\label{II.7.4.2}
We call an algebraic curve over $k$ any non-empty algebraic scheme over $k$
satisfying the conditions of (\hyperref[II.7.4.1]{7.4.1}).
\end{definition}

In the language of dimension, which will be introduced in Chapter IV, this
will be expressed by saying that an algebraic curve over $k$ is a non-empty
algebraic $k$-scheme \emph{all of whose irreducible components are of
dimension 1}.

Note that, if $X$ is an algebraic curve over $k$, the reduced closed
sub-preschemes $X_i$ $(1\leq i\leq n)$ of $X$ having as underlying spaces
the irreducible components of $X$ are algebraic curves over $k$.

\begin{corollary}[7.4.3]
\label{II.7.4.3}
Let $X$ be an irreducible algebraic curve. The only non-closed point of $X$ is
its generic point. The closed subsets of $X$ distinct from $X$ are the
finite sets of closed points; these are also the only subsets of $X$ that
are not everywhere dense.
\end{corollary}

If a point $x\in X$ is not closed, its closure in $X$ is an irreducible
closed subset of $X$, hence necessarily the whole of $X$ by virtue of
(\hyperref[II.7.4.1]{7.4.1}), and consequently $x$ is the generic point of
$X$. A closed subset $F$ of $X$ distinct from $X$ cannot contain

\oldpage[II]{149}
the generic point of $X$, hence all its points are closed (in $X$ and
\emph{a fortiori} in $F$); considering the reduced closed sub-prescheme of
$X$ having $F$ as underlying space (\hyperref[I.5.2.1]{I, 5.2.1}), it
therefore follows from (\hyperref[I.6.2.2]{I, 6.2.2}) that $F$ is finite and
discrete. The closure in $X$ of an infinite subset of $X$ is therefore
necessarily equal to $X$.

When $X$ is an arbitrary algebraic curve, by applying
(\hyperref[II.7.4.3]{7.4.3}) to the irreducible components of $X$, we see
that the only non-closed points of $X$ are the generic points of these
components.

\begin{corollary}[7.4.4]
\label{II.7.4.4}
Let $X$ and $Y$ be two irreducible algebraic curves over $k$ and $f:X\to Y$
a $k$-morphism. For $f$ to be dominant, it is necessary and sufficient that
$f^{-1}(y)$ be finite for every $y\in Y$.
\end{corollary}

Indeed, if $f$ is not dominant, $f(X)$ is necessarily a \emph{finite}
subset of $Y$ by virtue of (\hyperref[II.7.4.3]{7.4.3}), so it is not
possible for $f^{-1}(y)$ to be finite for every point of $Y$, since otherwise
$X$ would be finite, which is absurd (\hyperref[II.7.4.1]{7.4.1}).
Conversely, if $f$ is dominant, then, for every $y\in Y$ distinct from the
generic point $\eta$ of $Y$, $f^{-1}(y)$ is closed in $X$, since $\{y\}$ is
closed in $Y$ (\hyperref[II.7.4.3]{7.4.3}); on the other hand, by
hypothesis, $f^{-1}(y)$ does not contain the generic point $\xi$ of $X$, hence
is finite by virtue of (\hyperref[II.7.4.3]{7.4.3}). Finally, to see that
$f^{-1}(\eta)$ is finite when $f$ is dominant, we note that the fibre
$f^{-1}(\eta)$ is an irreducible scheme of finite type over $k(\eta)$ with
generic point $\xi$
(\hyperref[I.6.3.9]{I, 6.3.9} and
\hyperref[I.6.4.11]{6.4.11}). Since $k(\xi)$ and $k(\eta)$ are
extensions of finite type of $k$ with the same transcendence degree 1, $k(\xi)$
is necessarily an extension of finite degree of $k(\eta)$; hence $\xi$ is
closed in $f^{-1}(\eta)$ (\hyperref[I.6.4.2]{I, 6.4.2}), and $f^{-1}(\eta)$
consequently consists of the single point $\xi$.

We will see in Chapter III that a \emph{proper} morphism $f:X\to Y$ of
Noetherian preschemes such that $f^{-1}(y)$ is finite for every $y\in Y$ is
necessarily \emph{finite}; it will therefore follow from
(\hyperref[II.7.4.4]{7.4.4}) that a proper dominant morphism from an
irreducible algebraic curve to an algebraic curve is \emph{finite}.

\begin{corollary}[7.4.5]
\label{II.7.4.5}
Let $X$ be an algebraic curve over $k$. For $X$ to be regular, it is necessary
and sufficient that $X$ be normal, or again that the local rings of its
closed points be discrete valuation rings.
\end{corollary}

This follows immediately from condition b) of
(\hyperref[II.7.4.1]{7.4.1}) and from
(\hyperref[II.7.1.6]{7.1.6}).

\begin{corollary}[7.4.6]
\label{II.7.4.6}
Let $X$ be a reduced algebraic curve and $\mathcal{A}$ a reduced coherent
$\mathcal{R}(X)$-algebra; then the integral closure $X'$
of $X$ relative to $\mathcal{A}$ (\hyperref[II.6.3.4]{6.3.4}) is a
normal algebraic curve, and the canonical morphism $X'\to X$ is finite.
\end{corollary}

The fact that $X'\to X$ is finite follows from
(\hyperref[II.6.3.10]{6.3.10}); $X'$ is therefore an algebraic $k$-scheme;
moreover, if $x_i$ $(1\leq i\leq n)$ are the generic points of the irreducible
components of $X$ and $x'_j$ $(1\leq j\leq m)$ those of the irreducible
components of $X'$, each of the $k(x'_j)$ is a finite algebraic extension of
one of the $k(x_i)$ (\hyperref[II.6.3.6]{6.3.6}), hence has transcendence
degree 1 over $k$. $X'$ is therefore indeed an algebraic curve over $k$, and
moreover we know that $X'$ is the sum of finitely many integral and normal
schemes (\hyperref[II.6.3.6]{6.3.6} and
\hyperref[II.6.3.7]{6.3.7}).

\begin{env}[7.4.7]
\label{II.7.4.7}
We say that an algebraic curve $X$ over $k$ is \emph{complete} if it is
\emph{proper} over $k$.
\end{env}

\begin{corollary}[7.4.8]
\label{II.7.4.8}
For a reduced algebraic curve $X$ over $k$ to be complete, it is necessary and
sufficient that its normalization $X'$ be so.
\end{corollary}

\oldpage[II]{150}
Indeed, the canonical morphism $f:X'\to X$ is then finite
(\hyperref[II.7.4.6]{7.4.6}), hence proper
(\hyperref[II.6.1.11]{6.1.11}) and surjective
(\hyperref[II.6.3.8]{6.3.8}); if $g:X\to\operatorname{Spec}(k)$ is the
structural morphism, $g$ and $g\circ f$ are therefore simultaneously proper,
as follows from (\hyperref[II.5.4.2]{5.4.2, (iii)}) and
(\hyperref[II.5.4.3]{5.4.3, (ii)}), $g$ being separated by hypothesis.

\begin{proposition}[7.4.9]
\label{II.7.4.9}
Let $X$ be a normal algebraic curve over $k$ and $Y$ an algebraic $k$-scheme
proper over $k$. Then every rational $k$-map from $X$ to $Y$ is
everywhere defined; in other words, it is a morphism.
\end{proposition}

Indeed, it follows from (\hyperref[II.7.3.7]{7.3.7}) that, at the points
$x\in X$ where such a map is not defined, the dimension of $\mathcal{O}_x$
would have to be $\geq2$; hence the set of these points is empty; the last
assertion comes from (\hyperref[I.7.2.3]{I, 7.2.3}).

\begin{corollary}[7.4.10]
\label{II.7.4.10}
A normal algebraic curve over $k$ is quasi-projective over $k$.
\end{corollary}

Since $X$ is the sum of finitely many integral and normal algebraic curves
(\hyperref[II.6.3.8]{6.3.8}), we may restrict to the case where $X$ is
integral (\hyperref[II.5.3.6]{5.3.6}). Since $X$ is quasi-compact, it is
covered by finitely many affine open subsets $U_i$ $(1\leq i\leq n)$, and,
since each of these is of finite type over $k$, there exist an integer
$n_i$ and a $k$-immersion $f_i:U_i\to\mathbf{P}_k^{n_i}$
(\hyperref[II.5.3.3]{5.3.3} and
\hyperref[II.5.3.4]{5.3.4, (i)}). Since $U_i$ is dense in $X$, it follows
from (\hyperref[II.7.4.9]{7.4.9}) that $f_i$ extends to a $k$-morphism
$g_i:X\to\mathbf{P}_k^{n_i}$, whence a $k$-morphism
$g=(g_1,\ldots,g_n)_k$ from $X$ to the product $P$ of the
$\mathbf{P}_k^{n_i}$ over $k$. Moreover, for each index $i$, since the
restriction of $g_i$ to $U_i$ is an immersion, so is the restriction of $g$ to
$U_i$ (\hyperref[I.5.3.14]{I, 5.3.14}). Since the
$U_i$ cover $X$ and $g$ is separated
(\hyperref[I.5.5.1]{I, 5.5.1, (v)}), $g$ is an immersion of $X$ into $P$
(\hyperref[I.8.2.8]{I, 8.2.8}). Since the Segre morphism
(\hyperref[II.4.3.3]{4.3.3}) provides an immersion of $P$ into some
$\mathbf{P}_k^N$, this completes the proof that $X$ is quasi-projective.

\begin{corollary}[7.4.11]
\label{II.7.4.11}
A normal algebraic curve $X$ is isomorphic to the scheme induced on an
everywhere dense open subset of a normal and complete algebraic curve
$\widehat{X}$, determined up to unique isomorphism.
\end{corollary}

If $X_1$, $X_2$ are two normal and complete curves, it follows immediately
from (\hyperref[II.7.4.9]{7.4.9}) that every isomorphism of an open subset
$U_1$ dense in $X_1$ onto an open subset $U_2$ dense in $X_2$ extends uniquely
to an isomorphism of $X_1$ onto $X_2$; whence the uniqueness assertion. To
prove the existence of $\widehat{X}$, it suffices to remark that we can
regard $X$ as a subscheme of a projective bundle $\mathbf{P}_k^n$
(\hyperref[II.7.4.10]{7.4.10}). Let $\overline{X}$ be the closure of $X$ in
$\mathbf{P}_k^n$ (\hyperref[I.9.5.11]{I, 9.5.11}); since $X$ is induced
by $\overline{X}$ on an open subset dense in $\overline{X}$
(\hyperref[I.9.5.10]{I, 9.5.10}), the generic points $x_i$ of the
irreducible components of $X$ are also those of the irreducible components
of $\overline{X}$, and the $k(x_i)$ are the same for these two schemes; hence
(\hyperref[II.7.4.1]{7.4.1}) $\overline{X}$ is an algebraic curve over
$k$, which is reduced (\hyperref[I.9.5.9]{I, 9.5.9}) and projective over $k$
(\hyperref[II.5.5.1]{5.5.1}), hence complete
(\hyperref[II.5.5.3]{5.5.3}). Let us then take for $\widehat{X}$ the
\emph{normalization} of $\overline{X}$, which is again complete
(\hyperref[II.7.4.8]{7.4.8}); moreover, if
$h:\widehat{X}\to\overline{X}$ is the canonical morphism, the restriction of
$h$ to $h^{-1}(X)$ is an isomorphism onto $X$, since $X$ is normal
(\hyperref[II.6.3.4]{6.3.4}), and, since $h^{-1}(X)$ contains the generic
points of the irreducible components of $\widehat{X}$
(\hyperref[II.6.3.8]{6.3.8}), it is dense in $\widehat{X}$, which completes
the proof.

\oldpage[II]{151}

\begin{remark}[7.4.12]
\label{II.7.4.12}
We will show in Chapter V that the conclusion of
(\hyperref[II.7.4.10]{7.4.10}) is still valid without supposing the curve
normal (or even reduced); we will also show that an algebraic curve
(reduced or not) is affine if and only if none of its (reduced) irreducible
components is complete.
\end{remark}

\begin{corollary}[7.4.13]
\label{II.7.4.13}
Let $X$ be an irreducible normal curve with field $R(X)=K$ and $Y$ a complete
integral curve with field $L=R(Y)$. There is a canonical bijective
correspondence between the dominant $k$-morphisms $X\to Y$ and the
$k$-monomorphisms $L\to K$.
\end{corollary}

By virtue of (\hyperref[II.7.4.9]{7.4.9}), the rational $k$-maps
from $X$ to $Y$ are identified with the $k$-morphisms $u:X\to Y$. Since the
dominant morphisms $u$ are characterized by the fact that $u(x)=y$ (denoting by
$x$ and $y$ the respective generic points of $X$ and $Y$), the corollary
follows from these remarks and from (\hyperref[I.7.1.13]{I, 7.1.13}).

\begin{env}[7.4.14]
\label{II.7.4.14}
We can make the result of (\hyperref[II.7.4.13]{7.4.13}) more precise when we
take for $Y$ the \emph{projective line}
$\mathbf{P}_k^1=\operatorname{Proj}(k[T_0,T])$, $T_0$ and $T$ being two
indeterminates. This is indeed an integral scheme
(\hyperref[II.2.4.4]{2.4.4}), and the scheme induced on the open subset
$D_+(T_0)$ of $Y$ is isomorphic to $\operatorname{Spec}(k[T])$
(\hyperref[II.2.3.6]{2.3.6}); hence the generic point of $Y$ is the ideal
$(0)$ of $k[T]$ and its field of rational functions is $k(T)$, which shows
that $Y$ is a complete algebraic curve over $k$. Moreover, the only graded
prime ideal of $S=k[T_0,T]$ containing $T_0$ and distinct from $S_+$ is
the principal ideal $(T_0)$; hence the complement of $D_+(T_0)$ in
$Y=\mathbf{P}_k^1$ consists of \emph{a single closed point}, called the
``point at infinity'', which we will denote by $\infty$ (for a general study
of the relations between vector bundles and projective bundles, see
\hyperref[subsection:II.8.4]{8.4}). With this notation:
\end{env}

\begin{corollary}[7.4.15]
\label{II.7.4.15}
Let $X$ be an irreducible normal curve with field $R(X)=K$. There exists a
canonical bijective map of $K$ onto the set of morphisms $u$ from
$X$ to $\mathbf{P}_k^1$ distinct from the constant morphism with value
$\infty$. For $u$ to be dominant, it is necessary and sufficient that the
corresponding element of $K$ be transcendental over $k$.
\end{corollary}

This statement follows immediately from (\hyperref[II.7.4.9]{7.4.9}) and
from

\begin{lemma}[7.4.15.1]
\label{II.7.4.15.1}
Let $X$ be an integral prescheme over $k$, and let $K=R(X)$ be its field of
rational functions. There exists a canonical bijective map of the set
$K$ onto the set of rational maps $u$ from $X$ to
$\mathbf{P}_k^1$ distinct from the constant morphism with value $\infty$. For
such a rational map to be dominant, it is necessary and sufficient that the
corresponding element of $K$ be transcendental over $k$.
\end{lemma}

First of all, the rational maps from $X$ to $\mathbf{P}_k^1$
correspond bijectively to the points of $\mathbf{P}_k^1$ with values in
the extension $K$ of $k$ (\hyperref[I.7.1.12]{I, 7.1.12}). If such a point
is localized (\hyperref[I.3.4.5]{I, 3.4.5}) at the generic point of
$\mathbf{P}_k^1$, the corresponding rational map is evidently
dominant. Otherwise, since every point of $\mathbf{P}_k^1$
distinct from the generic point is closed (\hyperref[II.7.4.3]{7.4.3}),
the image of the domain of definition $U$ of $u$ under the unique morphism
$U\to\mathbf{P}_k^1$ of the class $u$ (\hyperref[I.7.2.2]{I, 7.2.2})
consists of a single closed point $y$ of $\mathbf{P}_k^1$, and this morphism
(which is not necessarily everywhere defined in $X$) is therefore not dominant;
by abuse of language, we then say that the rational map $u$ is
``constant, with value $y$''. It remains to put the points of
$\mathbf{P}_k^1$ with values in $K$ whose locality
(\hyperref[I.3.4.5]{I, 3.4.5}) is distinct from $\infty$ in bijective
correspondence with the elements of $K$, and to verify that the locality of
such a point is the generic point of $\mathbf{P}_k^1$ if and only if it
corresponds to an

\oldpage[II]{152}
element transcendental over $k$. Now this verification is immediate from
(\hyperref[II.4.2.6]{4.2.6, example 1\textsuperscript{o}}).

\begin{corollary}[7.4.16]
\label{II.7.4.16}
Let $X$, $Y$ be two normal, complete and irreducible algebraic curves over
$k$, and let $K=R(X)$, $L=R(Y)$ be their fields. There exists a
canonical bijective correspondence between the set of $k$-isomorphisms
$X\xrightarrow{\sim}Y$ and the set of $k$-isomorphisms
$L\xrightarrow{\sim}K$.
\end{corollary}

This is an evident consequence of (\hyperref[II.7.4.13]{7.4.13}).

\begin{env}[7.4.17]
\label{II.7.4.17}
Corollary (\hyperref[II.7.4.16]{7.4.16}) shows that a normal, complete
and irreducible algebraic curve over $k$ is \emph{determined by its field
of rational functions $K$ up to unique isomorphism}; by definition,
$K$ is an extension of finite type of $k$ of transcendence degree $1$ (what
is classically called an \emph{algebraic function field in one
variable}). Moreover:
\end{env}

\begin{proposition}[7.4.18]
\label{II.7.4.18}
For every extension $K$ of $k$ of finite type and of transcendence degree
$1$, there exists a normal, complete and irreducible algebraic curve $X$
such that $R(X)=K$ (determined up to unique isomorphism). The set of
local rings of $X$ is identified (\hyperref[I.8.2.1]{I, 8.2.1}) with
the set consisting of $K$ and of the valuation rings containing $k$ and having
$K$ as field of fractions.
\end{proposition}

Indeed, $K$ is an extension of finite degree of a purely transcendental
extension $k(T)$ of $k$, which is identified, as we have seen, with the field
of rational functions of the projective line $Y=\mathbf{P}_k^1$. Let $X$ be the
\emph{integral closure} of $Y$ relative to $K$
(\hyperref[II.6.3.4]{6.3.4}); $X$ is a normal algebraic curve with
field $K$ (\hyperref[II.6.3.7]{6.3.7}), and it is complete since the
morphism $X\to Y$ is finite (\hyperref[II.7.4.6]{7.4.6}). The local
rings $\mathcal{O}_x$ of $X$ are: the field $K$ when $x$ is the generic
point; if $x$ is distinct from the generic point, $\mathcal{O}_x$ is a
discrete valuation ring containing $k$ and having $K$ as field of
fractions (\hyperref[II.7.4.5]{7.4.5}). Conversely, let $A$ be such a
ring; since the morphism $X\to\operatorname{Spec}(k)$ is proper and
$A$ dominates $k$, it dominates a local ring $\mathcal{O}_x$ of $X$
(\hyperref[II.7.3.10]{7.3.10}); the latter, being a valuation ring
having $K$ as field of fractions, is therefore necessarily equal to $A$.

\begin{namedtheorem}[7.4.19]{Remarks}
\label{II.7.4.19}
It follows from (\hyperref[II.7.4.16]{7.4.16}) and
(\hyperref[II.7.4.18]{7.4.18}) that \emph{the datum of a normal, complete
and irreducible algebraic curve over $k$ is essentially equivalent to
the datum of an extension $K$ of $k$ of finite type and of transcendence
degree $1$}. Note that, if $k'$ is an extension of the base field
$k$, $X\otimes_k k'$ will again be a complete algebraic curve over $k'$
(\hyperref[II.5.4.2]{5.4.2, (iii)}), but in general it will be neither
reduced nor irreducible. It will be so, however, if $K$ is a separable
extension of $k$ and $k$ is algebraically closed in $K$ (which is expressed,
in a classical terminology that we will not follow, by saying that $K$ is a
``regular extension'' of $k$). But even in this case, it may happen that
$X\otimes_k k'$ is not normal. The reader will find details on these
questions in Chapter IV.
\end{namedtheorem}

